% !TeX encoding = UTF-8 !TeX spellcheck = en_US
\documentclass[11pt, a4paper]{article}

% The title of the current document to be produced.
\newcommand{\doctitle}{Assignment \#5}
\newcommand{\assignments}{\bluetext{\textbf{Assignment:}}}
%
\setlength{\unitlength}{1in}
\renewcommand{\arraystretch}{2}
\setlength{\columnsep}{2cm}

\newcommand\dunderline[3][-1pt]{{%
  \sbox0{#3}%
  \ooalign{\copy0\cr\rule[\dimexpr#1-#2\relax]{\wd0}{#2}}}}

\usepackage[sfdefault]{atkinson}
\usepackage[T1]{fontenc}

\def\code#1{\bluetext{\texttt{#1}}}


%------------------------------------------------------------
% Import commands for both teacher and course information.  | NOTE: Change your
% teacher and course info in these files. |
%------>------>------>------>------>------>------>------>-->|
%
%------------------------------------------------------------
%-- Import packages and custom command definitons.          |
%------>------>------>------>------>------>------>------>-->|
%% ==================================
%% ===== Teacher info             ===
%% ==================================
%%
\newcommand{\instructor}{Katelyn Breivik}
\newcommand{\office}{Wean 8319}
\newcommand{\hours}{3pm on Mondays}
\newcommand{\phone}{}
\newcommand{\college}{}
\newcommand{\email}{kbreivik@andrew.cmu.edu}
\newcommand{\faculty}{}
\newcommand{\department}{Department of Physics}
                              %|
%% ==================================
%% ===== Course-specific commands ===
%% ==================================

%- Instructions: change course info here. 
\newcommand{\semester}{Fall 2023}

\newcommand{\ponderation}{2-4-3 (Theory-Lab-Homework)}
\newcommand{\coursetitle}{Astrophysics of Stars and the Galaxy}
\newcommand{\coursenumber}{[33-467]}
\newcommand{\prerequisite}{33-224, 33-234}
 
\input{../includes/packages-imports}                          %|  
\input{../includes/custom-commands}   
%
%---> Genereate & inject metadata                           |
\input{../includes/hyperef.doc.info}                          %|
%------------------------------------------------------------

\topmargin -70pt
\begin{document} 

\normalfont

%-------------------------------------------------------------
%-- Make the header of the document                          |
%------>------>------>------>------>------>------>------>--> |
%---------------------------------------------------------------------
%-- The following produces the document header including the title.  |
%---------------------------------------------------------------------

\noindent % <-- need to have this first.

\hfill	
\begin{minipage}{0.60\textwidth}%
	\raggedleft%
	{\Large \textsc{\doctitle}\par}
	\doublerule 
\end{minipage}%
\vspace{0.9cm}

  
%-------------------------------------------------------------
%-- Problem # 1                                              |
%------>------>------>------>------>------>------>------>--> |
\noindent \dunderline{1.0pt}{\textbf{Problem \#1: Stellar Evolution (26
points)}}\\
We have officially covered all the top level things necessary to understand the
general principles of stellar evolution! This problem asks you to prove it.
\vspace{8pt}

\noindent \textbf{Problem \#1a: Accessing and plotting stellar models (10
pts)}\\
\noindent Since we don't have access to a supercomputer that can help us run many stellar
evolution models, we'll be using the MESA Isochrones and Stellar Tracks (MIST)
packaged grids. Your first task is to download the thoeretica evolutionary
tracks from the
\href{https://waps.cfa.harvard.edu/MIST/model_grids.html#eeps}{MIST website} for
nonrotating stars at solar metallicty ([Fe/H] = 0.0). 
\vspace{4pt}
\noindent Once you have the models, plot the evolution in the theoretical H-R Diagram of a
$1\,\rm{M}_{\odot}$, $5\,\rm{M}_{\odot}$, $10\,\rm{M}_{\odot}$, and
$20\,\rm{M}_{\odot}$ star on five separate figures. Be sure to label your axes
and make sure that hotter temperatures start at the left of the x-axis!
\vspace{8pt}

\noindent \textbf{Problem \#1b: Accessing and plotting stellar models (16
pts)}\\
\noindent For each star, label (either with text in the plotting software or by hand) the
critical phases of evolution that are included in each evolutionary track. You
should explain the reasoning for every slope and direction change that the star
makes as it evolves across the H-R Diagram for each of the stars you have
plotted. 


\vspace{10pt}
\noindent \dunderline{1.0pt}{\textbf{Problem \#2: Self Reflections in
Research (9 pts)}}\\
Now that you have some distance from your project 1 presentations, it's time to
reflect. Write down three things that worked well in your project and three
things that didn't. For each of the three things that didn't work, write down a
strategy for how you can mitigate running into those problems for project 2.

\end{document} 
